\clearpage
\SourceLeaf{105}%

\begin{center}
\large THEOREMA ELEGANTISSIMUM.
\end{center}

\bigskip

\noindent\hspace*{\parindent}Sit

\[ 1+\tfrac{1}{2}\cdot\tfrac{1}{2}xx+\tfrac{1}{2}\cdot\tfrac{1}{2}\cdot\tfrac{3}{4}\cdot\tfrac{3}{4}x^4+\tfrac{1}{2}\cdot\tfrac{1}{2}\cdot\tfrac{3}{4}\cdot\tfrac{3}{4}\cdot\tfrac{5}{6}\cdot\tfrac{5}{6}x^6+\text{etc.} = fx, \]

\[ \tfrac{1}{2}x+\tfrac{1}{2}\cdot\tfrac{3}{2}\cdot\tfrac{3}{4}x^3+\tfrac{1}{2}\cdot\tfrac{3}{2}\cdot\tfrac{3}{4}\cdot\tfrac{5}{4}\cdot\tfrac{5}{6}x^5+\text{etc.} = f'x, \]

\noindent eritque

\[ \sin\varphi\, f\sin\varphi\, f'\cos\varphi + \cos\varphi\, f\cos\varphi\, f'\sin\varphi = \frac{2}{\pi\sin\varphi\cos\varphi}. \]

\bigskip

\bigskip

\begin{center}
BEMERKUNGEN ZUM »THEOREMA ELEGANTISSIMUM«.
\end{center}

\noindent\hspace*{\parindent}Dieses Theorem hat \textsc{Gauss} auf die letzte Seite seines Handexemplars von »\textsc{Euler}, Methodus inveniendi curvas Maximi Minimive proprietate gaudentes« geschrieben. Der Satz fand hier Aufnahme wegen der eleganten Gestalt der Schlussformel, welche offenbar \textsc{Gauss} selbst frappirte. In der Sache handelt es sich um eine der einfachsten Differentialrelationen aus der Theorie des arithmetisch-geometrischen Mittels, wie sich mit Hülfe einiger weniger Formeln aus Art.\,13 und 14 des »Arithmetisch-geometrischen Mittels« (cf. ges. Werke III p.\,379 ff.) zeigen lässt. Es erweist sich nämlich $fx$ als die \textsc{Gauss}ische Reihe $F(\tfrac{1}{2},\tfrac{1}{2},1;x^2)$, so dass man auf Grund der beiden letzten Formeln p.\,381 l.\,c. setzen kann:

\[ x=\sin\varphi=\frac{b}{a}, \qquad \sqrt{1-x^2}=\cos\varphi=\frac{c}{a}, \]

\[ f\sin\varphi=\frac{a}{M( a,c)}, \qquad f\cos\varphi=\frac{a}{M( a,b)}. \]

\noindent Das hier aufgestellte Theorem kleidet sich daraufhin nach einer einfachen Zwischenrechnung in die Gestalt:

\[ \frac{b^2\mathord{.} d\log\dfrac{a}{M( a,b)}}{d\log\dfrac{c}{a}} - \frac{c^2\mathord{.} d\log\dfrac{a}{M( a,c)}}{d\log\dfrac{a}{b}} = \frac{2}{\pi}\,M( a,b)\,M( a,c), \]

\noindent eine Formel, die eine einfache Folge der ersten Gleichungen p.\,380 l.\,c. ist.

\hfill \textsc{Fricke}.

\bigskip

